\begin{abstract}
% We present PhyPrompt, a lightweight, model-agnostic front end that automatically refines text-to-video prompts for improved physical realism. Our approach employs a two-stage methodology.
% First, we fine-tune a large language model on a curated, physics-focused Chain-of-Thought dataset, enabling it to integrate core physics principles such as object motion and force interactions into user prompts while preserving their original intent.
% Second, we apply Group Relative Policy Optimization with a dynamic reward schedule that initially prioritizes semantic fidelity and then progressively shifts emphasis toward commonsense physics, driving the model to correct physical inconsistencies across diverse video generation models.
% Through zero-shot evaluation on the VideoPhy2 benchmark across  a few diffusion generators (Lavie, VideoCrafter2, CogVideoX), PhyPrompt demonstrates a remarkable 16.8\% improvement in combined semantic and physical success compared to baseline prompts.
% These results show that an RL-trained LLM can steer video synthesis toward adherence to real-world physical laws, without the need for manual prompt engineering.

% State-of-the-art text-to-video (T2V) generators frequently violate physical laws despite high visual quality. We demonstrate that this limitation stems not from model capability but from insufficient physical constraints in input prompts: manually adding physics details reliably produces physically plausible videos. However, manual prompt engineering requires expertise and does not scale.
% We present PhyPrompt, a reinforcement learning framework that automatically refines prompts to elicit physically realistic video generation. Our key innovation is a dynamic reward curriculum that resolves the fundamental conflict between semantic fidelity and physical realism. Remarkably, this curriculum achieves synergistic optimization: performance on both objectives simultaneously exceeds what either objective achieves when optimized alone (Semantic Adherence: 47.8\% vs. 45.6\% from SA-only; Physical Commonsense: 66.8\% vs. 65.2\% from PC-only). This demonstrates that our method discovers compositional prompt structures beyond conventional multi-objective trade-offs.
% PhyPrompt employs a lightweight 7B parameter model that outperforms GPT-4o (+3.8\% joint performance) and DeepSeek-V3 (+2.2\%, despite 100$\times$ more parameters), achieving 40.8\% combined semantic-physical success on VideoPhy2 benchmarks. Our approach transfers zero-shot across diverse T2V architectures (Lavie, VideoCrafter2, CogVideoX-5B), delivering up to 16.8\% improvement over baseline prompts without generator-specific fine-tuning. These results establish that domain-specialized reinforcement learning with compositional curricula can surpass general-purpose scaling for physics-aware generation.
% State-of-the-art text-to-video (T2V) generators frequently violate physical laws despite high visual quality. We show this stems from insufficient physical constraints in prompts rather than model limitations: manually adding physics details reliably produces physically plausible videos, but requires expertise and does not scale.
% We present PhyPrompt, a reinforcement learning framework that automatically refines prompts for physically realistic generation. Our dynamic reward curriculum achieves synergistic optimization: PhyPrompt-7B reaches 40.8\% joint success on VideoPhy2 (8.6 percentage points gain), improving physical commonsense by 11 percentage points (55.8\% to 66.8\%) while simultaneously increasing semantic adherence by 4.4pp (43.4\% to 47.8\%). Remarkably, our curriculum exceeds single-objective training on both metrics, demonstrating compositional prompt discovery beyond conventional multi-objective trade-offs.
% PhyPrompt outperforms GPT-4o (+3.8\% joint) and DeepSeek-V3 (+2.2\%, 100$\times$ larger) using only 7B parameters. The approach transfers zero-shot across diverse T2V architectures (Lavie, VideoCrafter2, CogVideoX-5B) with up to 16.8\% improvement, establishing that domain-specialized reinforcement learning with compositional curricula surpasses general-purpose scaling for physics-aware generation.
State-of-the-art text-to-video (T2V) generators frequently violate physical laws despite high visual quality. We show this stems from insufficient physical constraints in prompts rather than model limitations: manually adding physics details reliably produces physically plausible videos, but requires expertise and does not scale.
We present PhyPrompt, a two-stage reinforcement learning framework that automatically refines prompts for physically realistic generation. First, we fine-tune a large language model on a physics-focused Chain-of-Thought dataset to integrate principles like object motion and force interactions while preserving user intent. Second, we apply Group Relative Policy Optimization with a dynamic reward curriculum that initially prioritizes semantic fidelity, then progressively shifts toward physical commonsense. This curriculum achieves synergistic optimization: PhyPrompt-7B reaches 40.8\% joint success on VideoPhy2 (8.6pp gain), improving physical commonsense by 11pp (55.8\% to 66.8\%) while simultaneously increasing semantic adherence by 4.4pp (43.4\% to 47.8\%). Remarkably, our curriculum exceeds single-objective training on both metrics, demonstrating compositional prompt discovery beyond conventional multi-objective trade-offs.
PhyPrompt outperforms GPT-4o (+3.8\% joint) and DeepSeek-V3 (+2.2\%, 100$\times$ larger) using only 7B parameters. The approach transfers zero-shot across diverse T2V architectures (Lavie, VideoCrafter2, CogVideoX-5B) with up to 16.8\% improvement, establishing that domain-specialized reinforcement learning with compositional curricula surpasses general-purpose scaling for physics-aware generation.
\end{abstract}
